% Section 6 — Quantitative parity defect (charter §4.6; issues T1-5, T2-5)
% GO/NO-GO area (charter §2): OT/info-geometry survives only if it earns a theorem.
\section{Quantitative parity defect}
\label{sec:parity-defect}

\gonogo{Optimal transport + information geometry (T1-5 $\to$ T2-5) was
provisional. If the parity defect yielded only a restatement of the
information-theoretic content of \cref{sec:parity} and bought no new
proposition, this section was to be demoted to a remark and the decision
recorded in WOV-271 (T2-5). T1-5 supplied the foundations and \emph{defined}
the metric precisely enough to be computed against the R4 artifacts; the
explicit criterion is stated in \cref{subsec:defect-gonogo}.
\textbf{Resolved (T2-5, WOV-271): go.} The metric earns
\cref{thm:defect-risk-gap}: $\paritydefect$ upper-bounds the risk gap of a
\emph{fixed} bounded predictor across the
$\ladder{4}\!\leftrightarrow\!\ladderplus$ swap --- a quantity
\cref{thm:irrelevance} and the information gap $\infogap{\cdot}$ cannot see
(both vanish under parity) --- and \cref{cor:defect-finer-order} shows it
strictly refines the mutual-information ordering on the parity class, witnessed
by the R4 profile (\cref{ex:r4-defect}). The section stands
(\cref{subsec:defect-theorem}).}

% This file holds the T1-5 (WOV-266) FOUNDATIONS only: Wasserstein distance +
% Kantorovich-Rubinstein duality; statistical manifolds, the Fisher metric,
% KL as a local divergence; the DEFINITION of the quantitative parity defect
% (W1 between induced predictive distributions + a local Fisher sensitivity);
% and a recipe for computing it in principle against the R4 round-trip data.
% The metric is USED (and earns or loses its place) in T2-5.

\cref{sec:parity} treats parity as a binary equality of mutual informations,
and guardrail~G3 flags that real controls meet it only ``operationally, not
absolutely.'' This section turns ``how far from parity'' into a number. The
mutual-information gap $\MI{Y}{X}-\MI{Y}{\rep(X)}$ already measures lost
\emph{information}; what it does not measure is the residual asymmetry between
two representations that carry the \emph{same} information about $Y$ yet drive
a bounded predictor to slightly different predictive distributions --- exactly
the $\sigalg{\ladder{4}}\ne\sigalg{\ladderplus}$ gap of
\cref{prop:abs-strict}. We build the optimal-transport and
information-geometry tools that quantify that residual, define the parity
defect $\paritydefect$, and specify how it is computed against the R4
round-trip artifacts. We follow \citet{villani2009} and \citet{peyrecuturi2019}
for optimal transport and \citet{amarinagaoka2000} for information geometry.

\subsection{Wasserstein distance and Kantorovich--Rubinstein duality}
\label{subsec:wasserstein}

\begin{definition}[Wasserstein distance]
\label{def:wasserstein}
Let $(\Aspace,d)$ be a Polish metric space and $\mu,\nu$ probability measures
on it with finite first moment. The \emph{$1$-Wasserstein distance} is
\[
  \Wass{1}(\mu,\nu)
  \;\defeq\;
  \inf_{\gamma\in\Pi(\mu,\nu)}\;
  \int_{\Aspace\times\Aspace} d(a,a')\,\mathrm{d}\gamma(a,a'),
\]
the infimum over couplings $\gamma$ with marginals $\mu$ and $\nu$. It is the
minimal average cost of transporting $\mu$ onto $\nu$.
\end{definition}

\begin{proposition}[Kantorovich--Rubinstein duality]
\label{prop:kr-duality}
With $\mu,\nu$ as above,
\begin{equation}
  \label{eq:kr}
  \Wass{1}(\mu,\nu)
  \;=\;
  \sup_{\,\|g\|_{\mathrm{Lip}}\le 1}\;
  \Big(\,\Expect_{\mu}[g] - \Expect_{\nu}[g]\,\Big),
\end{equation}
the supremum over $1$-Lipschitz test functions $g:\Aspace\to\mathbb{R}$.
\end{proposition}

\begin{proof}
This is the dual form of the Kantorovich problem for the cost $d$; see
\citet[Thm.~5.10 and Particular Case~5.16]{villani2009}.
\end{proof}

The duality \eqref{eq:kr} is what makes the defect \emph{estimable}: it
rewrites a transport cost over couplings as an expectation gap over Lipschitz
witnesses, which can be evaluated from samples of the two predictive
distributions without solving the primal coupling problem.

\subsection{Statistical manifolds and the Fisher metric}
\label{subsec:fisher}

Where the Wasserstein distance measures the global cost of moving one
predictive distribution to another, the Fisher metric measures the
\emph{local} sensitivity of a predictor to small representation
perturbations, which is the regime relevant when $\ladderplus$ and
$\ladder{4}$ are close.

\begin{definition}[Statistical manifold, Fisher metric]
\label{def:fisher}
A parametric family $\{p_\theta:\theta\in\Theta\}$ of densities, smooth in
$\theta$ over an open $\Theta\subseteq\mathbb{R}^k$, is a \emph{statistical
manifold}. Its \emph{Fisher information metric} at $\theta$ is the
$k\times k$ matrix
\[
  \Fisher(\theta)_{ij}
  \;\defeq\;
  \Expect_{p_\theta}\!\Big[
    \partial_{\theta_i}\log p_\theta \;
    \partial_{\theta_j}\log p_\theta
  \Big],
\]
the covariance of the score. It is the unique (up to scale) Riemannian metric
invariant under sufficient statistics \citep[Ch.~2--3]{amarinagaoka2000}.
\end{definition}

\begin{proposition}[KL is the Fisher metric to second order]
\label{prop:kl-fisher}
For $\theta'=\theta+\mathrm{d}\theta$,
\[
  \KL{p_\theta}{p_{\theta'}}
  \;=\;
  \tfrac12\,\mathrm{d}\theta^\top \Fisher(\theta)\,\mathrm{d}\theta
  \;+\; o(\|\mathrm{d}\theta\|^2).
\]
\end{proposition}

\begin{proof}
Taylor-expand $\theta'\mapsto\KL{p_\theta}{p_{\theta'}}$ at $\theta'=\theta$:
the value and gradient vanish, and the Hessian is $\Fisher(\theta)$
\citep[Ch.~3]{amarinagaoka2000}.
\end{proof}

\cref{prop:kl-fisher} ties this section back to \cref{sec:parity}: the KL
divergence that defines mutual information (\cref{def:mi}) is, locally, the
Fisher quadratic form. So a Fisher-sensitivity term and a relative-entropy
term are the same object at two scales --- which is why the defect below pairs
a global Wasserstein term with a local Fisher term rather than choosing
between them.

\subsection{The quantitative parity defect}
\label{subsec:defect-def}

Let $\preddist{c}$ denote the \emph{induced predictive distribution} under
condition $c$: the law on the action space $\Aspace$ of the bounded
predictor's output $\kernel{K}(\rep_c(X),\cdot)$ when $X\sim\Prob_X$, for
$c\in\{\ladder{4},\ladderplus\}$. Under \emph{exact} parity these two laws
coincide for the Bayes predictor; for a bounded predictor they may differ even
when $\MI{Y}{\ladder{4}(X)}=\MI{Y}{\ladderplus(X)}$.

\begin{definition}[Quantitative parity defect]
\label{def:parity-defect}
Fix a bounded predictor $\kernel{K}$ and a perturbation scale. The
\emph{parity defect} between $\ladder{4}$ and $\ladderplus$ is
\begin{equation}
  \label{eq:parity-defect}
  \paritydefect
  \;\defeq\;
  \underbrace{\Wass{1}\!\big(\preddist{\ladder{4}},\,\preddist{\ladderplus}\big)}_{\text{global: predictive-law transport}}
  \;+\;
  \lambda\,
  \underbrace{\sqrt{\,\overline{v}^\top \Fisher\, \overline{v}\,}}_{\text{local: Fisher sensitivity}},
\end{equation}
where $\overline{v}$ is the mean representation perturbation induced by the
round trip $\roundtrip{\ladder{4}}{\ladderplus}$ (\cref{def:cat-parity}),
identified by \cref{cor:defect-subadditive} with the complement of the
idempotent splitting of \cref{thm:retraction},
$\Fisher$ is the Fisher metric of the predictor's output family at the
operating point, and $\lambda\ge0$ trades off the two scales. Exact parity is
$\paritydefect=0$; $\paritydefect>0$ certifies operational-but-not-absolute
parity and bounds, via \cref{prop:kr-duality}, the largest difference in
expected Lipschitz functionals of the two predictive distributions.
\end{definition}

The two terms are complementary by \cref{prop:kl-fisher}: the Wasserstein term
captures the part of the asymmetry that survives at the global scale of the
output distribution, and the Fisher term captures the local sensitivity of the
predictor to the residual structural perturbation that the round trip fails to
undo. \cref{eq:parity-defect} is a definition, not yet a theorem; T2-5 must
show $\paritydefect$ either bounds a risk gap or constrains the predictive
asymmetry in a way the mutual-information gap of \cref{sec:parity} does not
(\cref{subsec:defect-gonogo}).

\subsection{Computing the defect against the R4 round-trip data}
\label{subsec:defect-r4}

The R4 round-trip recovery experiment (empirical paper Appendix~D) re-parses
$\ladderplus$ back to a Contract IR $\ladder{4}'$ and reports how much of
$\ladder{4}$'s structure is recovered: contracts $100\%$, AST node kinds
$99.8\%$, types $96.4\%$, parent--child edges $98.7\%$, byte-exact AST
$0/30$. This profile is, in the categorical reading of
\cref{subsec:retraction-parity}, the empirical signature of the idempotent
$e\circ p$: near-identity on semantic structure, far from identity on surface
bytes. The defect \eqref{eq:parity-defect} is computed against it as follows.

\begin{enumerate}
  \item \textbf{Perturbation $\overline{v}$ (Fisher term).} Each
    non-recovered structural element of the round trip is a coordinate of the
    representation perturbation; the recovery rates above estimate the
    per-coordinate failure probabilities, and $\overline{v}$ is their
    aggregate over the $30$ functions. Byte-level coordinates (recovery $0/30$)
    dominate $\overline{v}$ in magnitude but are predictively inert, which is
    why the Fisher term --- weighting perturbations by their effect on the
    \emph{output} family --- and not the raw byte distance is the right local
    measure.
  \item \textbf{Predictive laws (Wasserstein term).} Sample the bounded
    predictor on the C4 inputs and on their re-parsed $\ladder{4}'$
    counterparts to obtain empirical $\preddist{\ladder{4}}$ and
    $\preddist{\ladderplus}$, then estimate
    $\Wass{1}(\preddist{\ladder{4}},\preddist{\ladderplus})$ either by the
    primal optimal-transport solver of \citet[§2--3]{peyrecuturi2019} or by
    maximizing the Kantorovich--Rubinstein dual \eqref{eq:kr} over a Lipschitz
    witness class.
  \item \textbf{Trade-off $\lambda$.} Fixed by matching the two terms' units on
    a held-out calibration slice; the procedure is specified with the T2-5
    computation, not chosen post hoc.
\end{enumerate}

\begin{remark}[Status of the computation, G3/G6]
\label{rem:defect-status}
The recipe is ``computable in principle'': it specifies estimable quantities
and named estimators, but it is \emph{not} executed in this paper, and no
numerical value of $\paritydefect$ is claimed. The R4 rates are used as the
inputs the metric \emph{would} consume, illustrating that the defect is
grounded in measured artifacts rather than postulated. Any statement that the
computed defect is small (and hence that the empirical control was close to
exact parity) is, until T2-5 runs it, a conjecture and is labeled as such
(guardrail~G6).
\end{remark}

\subsection{Go/no-go criterion handed to T2-5}
\label{subsec:defect-gonogo}

\begin{quote}
\textbf{Criterion (depth of OT/info-geometry).} Section~\ref{sec:parity-defect}
survives as a section iff T2-5 establishes that $\paritydefect$
\eqref{eq:parity-defect} (i) upper-bounds the bounded predictor's risk gap
between $\ladder{4}$ and $\ladderplus$, or (ii) provides a strictly finer
ordering of representations than the mutual-information gap of
\cref{sec:parity} (two reps with equal information gap but different defect,
witnessed against the R4 profile). If T2-5 produces only a reformulation of
the information gap, the section is demoted to a remark recording that the
parity defect coincides with the information gap in this setting, and the
decision is logged in WOV-271.
\end{quote}

\subsection{What the metric earns: the defect bounds the bounded-predictor risk gap (T2-5)}
\label{subsec:defect-theorem}

This subsection discharges the go/no-go of \cref{subsec:defect-gonogo}, and the
verdict is \textbf{go}. The criterion asks for \emph{either} prong; the defect
delivers both. We show (i) $\paritydefect$ upper-bounds the risk gap of a fixed
bounded predictor across the $\ladder{4}\leftrightarrow\ladderplus$ swap
(\cref{thm:defect-risk-gap}, with the Fisher term supplying the local regime in
\cref{prop:defect-fisher-local}), and (ii) it strictly refines the
mutual-information ordering on the parity class
(\cref{cor:defect-finer-order}), a separation realized by the R4 control pair
(\cref{ex:r4-defect}). Full proofs are in \cref{app:parity-defect}; the decision
is recorded in WOV-271.

The object both prongs concern is the \emph{shared reader}: a single bounded
predictor consumes either representation. Formally we keep exactly the setup
\cref{def:parity-defect} already commits to --- one kernel $\kernel{K}$ evaluated
at $\rep_c(X)$ for $c\in\{\ladder{4},\ladderplus\}$ --- which models the frozen
LLM of \texttt{preprint-v1} reading $\ladder{4}$ or $\ladderplus$ through their
common token-sequence codomain. Its risk on condition $c$ is
$\risk(\kernel{K};c)=\Expect_{(X,Y)}\!\int_{\Aspace}\loss(Y,a)\,\kernel{K}(\rep_c(X),\mathrm{d}a)$
(\cref{def:risk}), and the quantity at issue is the
\emph{bounded-predictor representation gap}
\begin{equation}
  \label{eq:bp-gap}
  \advantage_{\kernel{K}}(\ladder{4},\ladderplus)
  \;\defeq\;
  \risk(\kernel{K};\ladder{4})-\risk(\kernel{K};\ladderplus).
\end{equation}
This is \emph{not} the advantage $\advantage(\rep,\rep')$ of \cref{def:advantage}:
that compares two \emph{independently optimized} ERMs and is capacity-bounded
(\cref{thm:bounded-gap}); \eqref{eq:bp-gap} compares \emph{one fixed} predictor
across the two formats and is, as we now show, transport-bounded. The two are
different operational quantities, which is exactly why \cref{sec:parity-defect}
is not subsumed by \cref{sec:bounded}.

\begin{definition}[Input-coupled predictive transport]
\label{def:coupled-transport}
For the shared predictor $\kernel{K}$ define
\begin{equation}
  \label{eq:coupled-transport}
  \cpredtransport
  \;\defeq\;
  \Expect_X\!\Big[\,
    \Wass{1}\big(\kernel{K}(\ladder{4}(X),\cdot),\,\kernel{K}(\ladderplus(X),\cdot)\big)
  \,\Big],
\end{equation}
the input-coupled form of the global term of \eqref{eq:parity-defect}. By joint
convexity of $\Wass{1}$ \citep[Thm.~4.8]{villani2009} and
$\preddist{c}=\Expect_X[\kernel{K}(\rep_c(X),\cdot)]$,
\begin{equation}
  \label{eq:marginal-le-coupled}
  \Wass{1}\big(\preddist{\ladder{4}},\preddist{\ladderplus}\big)\;\le\;\cpredtransport,
\end{equation}
with equality when the input pairing already realizes an optimal coupling of the
marginals (e.g.\ a swap-insensitive predictor). The marginal term of
\cref{def:parity-defect} is thus the relaxation; $\cpredtransport$ is the
risk-controlling form, and it inherits the rung-wise subadditivity of
\cref{cor:defect-subadditive} ($\Expect_X\Wass{1}$ is again a metric on the
predictive laws).
\end{definition}

\begin{theorem}[The defect dominates the bounded-predictor risk gap]
\label{thm:defect-risk-gap}
Let $\kernel{K}$ be a fixed bounded predictor read on both representations and
suppose the per-instance loss $a\mapsto\loss(y,a)$ is $L$-Lipschitz on
$(\Aspace,d)$ uniformly in $y$. Then
\begin{equation}
  \label{eq:defect-risk-gap}
  \big|\advantage_{\kernel{K}}(\ladder{4},\ladderplus)\big|
  \;\le\;
  L\,\cpredtransport,
\end{equation}
the global (Wasserstein) term of the parity defect \eqref{eq:parity-defect}.
Under operational parity $\ladder{4}\parityeq\ladderplus$ the information gap
contributes nothing --- $\infogap{\ladder{4}}=\infogap{\ladderplus}=0$
(\cref{prop:suff-iff-mi}), so the middle term of the excess-risk decomposition
\eqref{eq:excess} cancels between the two conditions --- and the entire surviving
gap of the shared predictor is controlled by the transport defect, with no
mutual-information term.
\idealization{This relaxes idealization~(i) of \cref{rem:slogan}
(Bayes-optimality): $\kernel{K}$ is an arbitrary fixed bounded predictor, not the
posterior $\post{\rep}$. It engages the failure of idealization~(ii) in its sharp
form --- the bound is informative precisely when
$\sigalg{\ladder{4}}\ne\sigalg{\ladderplus}$ (\cref{prop:abs-strict}) makes
$\kernel{K}$ respond to the swap; under absolute parity both sides are $0$. The
bound limits the \emph{size} of the swap gap in either direction and asserts no
effect exists (guardrail~G2); the Lipschitz-loss hypothesis is discharged in the
near regime by \cref{prop:defect-fisher-local} for merely bounded losses.}
\end{theorem}

\begin{proof}[Proof idea]
For fixed $(x,y)$, $\loss(y,\cdot)/L$ is $1$-Lipschitz, so
Kantorovich--Rubinstein duality \eqref{eq:kr} bounds the per-instance gap of
expected loss by $L\,\Wass{1}(\kernel{K}(\ladder{4}(x),\cdot),
\kernel{K}(\ladderplus(x),\cdot))$, uniformly in $y$. Take $\Expect_{(X,Y)}$ and
pull the absolute value inside. The information-gap cancellation is
\eqref{eq:excess} read for the shared kernel under \cref{prop:suff-iff-mi}. Full
proof in \cref{app:parity-defect}.
\end{proof}

The Lipschitz hypothesis is the right one when the two predictive laws are far
apart, but it says nothing when the loss is merely bounded (the regime of
\cref{cor:ratedist-gap}(b)). There the \emph{local} Fisher term takes over,
exactly as \cref{prop:kl-fisher} anticipated.

\begin{proposition}[Local Fisher control of the swap gap]
\label{prop:defect-fisher-local}
Suppose instead $\loss$ is bounded, $\|\loss\|_\infty<\infty$, and the predictor's
output family is smooth in the representation perturbation, the round trip
$\roundtrip{\ladder{4}}{\ladderplus}$ inducing the mean perturbation
$\overline{v}$ of \cref{def:parity-defect}. Then
\begin{equation}
  \label{eq:defect-fisher-local}
  \big|\advantage_{\kernel{K}}(\ladder{4},\ladderplus)\big|
  \;\le\;
  \|\loss\|_\infty\,\sqrt{\tfrac14\,\overline{v}^\top \Fisher\, \overline{v}}
  \;+\; o(\|\overline{v}\|),
\end{equation}
so to leading order the swap gap is governed by the Fisher term of
\eqref{eq:parity-defect}.
\idealization{Like \cref{thm:defect-risk-gap} this is a bounded-predictor
statement (idealization~(i) relaxed). It is a local (small-$\overline{v}$)
bound: the $o(\|\overline{v}\|)$ remainder is the Taylor error of
\cref{prop:kl-fisher}, and the inequality is informative only when the round-trip
perturbation is small --- the near-parity regime in which the global bound
\eqref{eq:defect-risk-gap} may be loose. Sign-agnostic (G2).}
\end{proposition}

\begin{proof}[Proof idea]
Per instance, the bounded-loss/total-variation step of
\cref{cor:ratedist-gap}(b) gives
$|\!\int\loss(y,\cdot)\,\mathrm{d}(p-p')|\le\|\loss\|_\infty\|p-p'\|_{\mathrm{TV}}$;
Pinsker bounds $\|p-p'\|_{\mathrm{TV}}\le\sqrt{\tfrac12\KL{p}{p'}}$ and
\cref{prop:kl-fisher} expands
$\KL{p}{p'}=\tfrac12\overline{v}^\top\Fisher\overline{v}+o(\|\overline{v}\|^2)$ at
the aggregate increment $\overline{v}$ of \cref{def:parity-defect}. The bound is
then $x$-free, so $\Expect_{(X,Y)}$ leaves it unchanged. Full proof in
\cref{app:parity-defect}.
\end{proof}

\begin{corollary}[The two terms are complementary, not redundant]
\label{cor:defect-two-term}
With $\lambda$ fixed by the unit-matching calibration of \cref{subsec:defect-r4},
there is a constant $c=c(L,\|\loss\|_\infty,\lambda)$ with
\[
  \big|\advantage_{\kernel{K}}(\ladder{4},\ladderplus)\big|
  \;\le\; c\,\paritydefect \;+\; o(\|\overline{v}\|),
\]
the global term \eqref{eq:defect-risk-gap} dominating the far regime (Lipschitz
loss, predictive laws far apart) and the Fisher term
\eqref{eq:defect-fisher-local} the near regime (bounded loss, small
perturbation). Each term bounds the swap gap exactly where the other is loose or
unavailable, which is why \eqref{eq:parity-defect} carries both rather than
choosing between them.
\end{corollary}

\begin{proof}
Take the smaller of the two right-hand sides of \eqref{eq:defect-risk-gap} and
\eqref{eq:defect-fisher-local} and bound it by their sum, absorbing
$L,\|\loss\|_\infty,\lambda$ into $c$.
\end{proof}

This settles prong~(i). Prong~(ii) is now immediate and is what makes the defect
strictly more than a re-encoding of \cref{sec:parity}.

\begin{corollary}[The defect strictly refines the information-gap ordering]
\label{cor:defect-finer-order}
On the class of representations in operational parity with $\ladder{4}$, the
information gap $\infogap{\cdot}$ is constant ($\equiv 0$,
\cref{prop:suff-iff-mi}) and so, by \cref{thm:irrelevance}, is the Bayes risk:
neither orders the class. The parity defect does. It is generically positive
there --- $\paritydefect>0$ whenever $\sigalg{\ladder{4}}\ne\sigalg{\ladderplus}$
(\cref{prop:abs-strict}) and $\kernel{K}$ is swap-sensitive, the case
\cref{thm:retraction}(iii) places in the common kernel of $\infogap{\cdot}$ and
$\bayesrisk(\cdot)$ --- and by \cref{thm:defect-risk-gap} it ranks parity pairs by
the bounded-predictor risk gap they admit. Hence $\paritydefect$ separates
representations the information gap and the Bayes risk cannot, a strict
refinement of both.
\end{corollary}

\begin{proof}
Constancy of $\infogap{\cdot}$ and $\bayesrisk(\cdot)$ on the parity class is
\cref{prop:suff-iff-mi} and \cref{thm:irrelevance}. Positivity of
$\paritydefect$ off absolute parity is \cref{def:parity-defect} with
\cref{prop:abs-strict}; that it bounds (hence orders by) the swap gap is
\cref{thm:defect-risk-gap}. The control pair
$(\ladder{4},\ladderplus)$ has $\infogap{\ladder{4}}=\infogap{\ladderplus}=0$ by
design yet $\paritydefect>0$ by \cref{ex:r4-defect}, an explicit witness.
\end{proof}

\begin{example}[The defect on the R4 round trip: an illustrative computation]
\label{ex:r4-defect}
We instantiate \eqref{eq:parity-defect} on the R4 profile by the recipe of
\cref{subsec:defect-r4}. Each non-recovered structural element is a coordinate of
$\overline{v}$ with magnitude its round-trip failure rate ($1-$recovery);
\cref{tab:r4-defect} collects them. The Fisher term weights each coordinate by
its effect on the predictor's output: the byte-exact coordinate, recovered
$0/30$ and hence maximal in raw magnitude, is \emph{predictively inert} and
carries Fisher weight $\approx 0$, while the semantic coordinates carry unit
weight (an illustrative diagonal $\Fisher$). With the action diameter normalized
to $1$ and $L=\lambda=1$, the global term is bounded by the predictively-relevant
failure mass and the local term by the Fisher quadratic form:
\[
  \cpredtransport \;\lesssim\; 0.051,
  \qquad
  \sqrt{\overline{v}^\top\Fisher\,\overline{v}} \;=\; 0.038,
  \qquad
  \paritydefect \;\approx\; 0.089 .
\]
The contrast is the point. A naive defect that used raw structural distance ---
unit weight on \emph{every} coordinate, byte coordinate included --- would read
$\|\overline{v}\|\approx 1.00$, larger by a factor of $\approx 26$ and dominated
entirely by the byte-recovery catastrophe. The Fisher weighting is exactly what
collapses that $1.00$ to $0.04$, certifying the control as \emph{near} parity
despite $0/30$ byte recovery. This is work neither the raw byte distance (which
sees a near-maximal perturbation) nor the information gap (which is identically
$0$) can do, and it is the quantitative form of the qualitative statement of
\cref{subsec:retraction-parity} that the round-trip idempotent is near-identity
on semantic structure and far from it on surface bytes.
\end{example}

\begin{table}[t]
\centering
\small
\begin{tabular}{@{}lrrl@{}}
\toprule
Round-trip coordinate & Recovery & $\overline{v}_i$ ($=1-$rec.) & Predictive role \\
\midrule
Contracts            & $100\%$  & $0.000$ & relevant (weight $1$) \\
AST node kinds       & $99.8\%$ & $0.002$ & relevant (weight $1$) \\
Parent--child edges  & $98.7\%$ & $0.013$ & relevant (weight $1$) \\
Types                & $96.4\%$ & $0.036$ & relevant (weight $1$) \\
Byte-exact AST       & $0/30$   & $1.000$ & \emph{inert} (weight $\approx 0$) \\
\midrule
\multicolumn{2}{@{}l}{Global term $\cpredtransport\lesssim\sum_{\text{rel.}}\overline{v}_i$} & $0.051$ & union-bound proxy \\
\multicolumn{2}{@{}l}{Fisher term $\sqrt{\overline{v}^\top\Fisher\overline{v}}$} & $0.038$ & semantic coords only \\
\multicolumn{2}{@{}l}{Illustrative $\paritydefect$ ($L=\lambda=1$)} & $0.089$ & \eqref{eq:parity-defect} \\
\bottomrule
\end{tabular}
\caption{The parity defect computed on the R4 round-trip recovery profile
(\cref{subsec:defect-r4}, \cref{ex:r4-defect}). Numbers are \emph{illustrative}:
action diameter normalized to $1$, $L=\lambda=1$, and a diagonal Fisher metric
with unit weight on the predictively-relevant coordinates and zero on the inert
byte coordinate. They are not the output of running the frozen predictor
(guardrail~G6; \cref{rem:defect-status}). The naive raw-distance defect
$\|\overline{v}\|\approx 1.00$ is $\approx 26\times$ larger and byte-dominated;
the Fisher weighting is what makes the defect small.}
\label{tab:r4-defect}
\end{table}

\begin{remark}[Why this passes the criterion; scope of the go (G4/G6)]
\label{rem:defect-go-scope}
\cref{thm:defect-risk-gap} and \cref{cor:defect-finer-order} meet the criterion of
\cref{subsec:defect-gonogo} on \emph{both} prongs (it required either). The defect
is not a reformulation of the information gap because the quantity it controls ---
the swap risk gap \eqref{eq:bp-gap} of a fixed bounded predictor --- is provably
in the common kernel of $\infogap{\cdot}$ and $\bayesrisk(\cdot)$
(\cref{thm:retraction}(iii)): both are $0$ across the parity pair while the defect
is positive (\cref{ex:r4-defect}). So a proposition that bounds and orders by the
swap gap cannot be a restatement of either. The two terms are shown
non-redundant by \cref{cor:defect-two-term}, and the R4 example shows the Fisher
term in particular does work the raw byte distance cannot. \textbf{Scope (G6).}
The numeric defect of \cref{ex:r4-defect} is illustrative and computed-in-principle
(\cref{rem:defect-status}): it consumes the measured R4 rates but normalizes the
action geometry, fixes $L,\lambda$, and posits a diagonal Fisher metric rather
than estimating one from the frozen predictor's outputs. The go rests on the
propositions \cref{thm:defect-risk-gap}, \cref{prop:defect-fisher-local},
\cref{cor:defect-finer-order}; the R4 example anchors them in artifacts and
exhibits the prong-(ii) separation, and no claim is made that the executed value
is small until T2-5's recipe is actually run (a conjecture, so labeled).
\end{remark}
